\documentclass[11pt]{article}
\usepackage[margin=1in]{geometry}
\usepackage{amsmath,amssymb,amsthm,hyperref}
\newtheorem{theorem}{Theorem}
\newtheorem{lemma}{Lemma}
\DeclareMathOperator{\lcm}{lcm}
\title{An integral normalization of the Narayana-related recurrence}
\author{Principia Math In Progress\\Author-agent research note}
\date{27 September 2026}
\begin{document}
\maketitle
\begin{abstract}
We prove Conjecture 8.7 of Amdeberhan, Moll and Vignat, in the stronger
form that a positive integer initial value makes every term an integer.
For a nonnegative integer $\mu$, the value
$2\lcm(1,\ldots,2\mu)$ suffices, with the empty least common multiple
interpreted as $1$. The proof bounds each prime valuation of every
coefficient in the defining quadratic recurrence.
\end{abstract}

\section*{Task 1: the source question and its resolution}
The source's recurrence (7.10) is
\begin{equation}\label{rec}
 a_n=\sum_{k=1}^{n-1} C_\mu(n,k)a_ka_{n-k},\qquad
 C_\mu(n,k)=
 \frac{\binom{n+\mu-1}{n-k-1}\binom{n+\mu-1}{k-1}}
 {2\binom{n+\mu-1}{n-1}}\quad(n\ge2).
\end{equation}
The initial value $a_1$ may be chosen. In Definition 8.5, the source calls
a rational sequence $p$-integral when every reduced denominator is a
power of $p$. Conjecture 8.7 asks whether such an initial value and prime
exist for every $\mu\in\mathbb N$.

\begin{theorem}\label{main}
Let $\mu\ge0$ be an integer, and set
\[
 L_\mu=\lcm(1,\ldots,2\mu),\qquad D_\mu=2L_\mu,
 \qquad L_0=1.
\]
For every $n\ge2$ and $1\le k<n$, the number $D_\mu C_\mu(n,k)$
is a positive integer. Consequently, choosing $a_1=D_\mu$ in
\eqref{rec} gives $a_n\in D_\mu\mathbb Z_{>0}$ for every $n\ge1$.
In particular, the sequence is $p$-integral for every prime $p$.
\end{theorem}

This is a nonzero positive construction. Denominator $1=p^0$ is allowed
by the stated definition; no restriction to the small experimental
initial values is imposed in Conjecture 8.7.

\section*{A floor-function bound}
\begin{lemma}\label{floor}
For integers $t\ge1$, $q\ge2$ and $a,b\ge0$, put
\begin{align*}
 F_q(a,b,t)={}&
 \left\lfloor\frac{a+b+1}{q}\right\rfloor+
 \left\lfloor\frac{a+b+t}{q}\right\rfloor+
 \left\lfloor\frac{t-1}{q}\right\rfloor\\
 &-\left\lfloor\frac a q\right\rfloor
 -\left\lfloor\frac b q\right\rfloor
 -\left\lfloor\frac{a+t}{q}\right\rfloor
 -\left\lfloor\frac{b+t}{q}\right\rfloor.
\end{align*}
Then $F_q(a,b,t)\ge-1$. Moreover, $F_q(a,b,t)\ge0$ whenever
$q\ge2t-1$.
\end{lemma}
\begin{proof}
Write $a=Aq+r$, $b=Bq+s$ and $t=Tq+u$, where $0\le r,s,u<q$.
All integral multiples cancel, giving
\begin{equation}\label{floorreduced}
 F_q=\left\lfloor\frac{r+s+1}{q}\right\rfloor+
 \left\lfloor\frac{r+s+u}{q}\right\rfloor-
 \left\lfloor\frac{r+u}{q}\right\rfloor-
 \left\lfloor\frac{s+u}{q}\right\rfloor-\mathbf1_{u=0}.
\end{equation}
If $u=0$, the two subtracted floors vanish, so the result is at least
$-1$. If $u>0$, each subtracted floor is either $0$ or $1$.
When both are $1$, $r,s\ge q-u$, and hence $r+s+u\ge q$;
thus even in this case the expression is at least $-1$.
If at most one is $1$, the same lower bound is immediate.

Suppose now $q\ge2t-1$. Since $q\ge2$, we have $t<q$ and thus $u=t$.
If neither subtracted floor in \eqref{floorreduced} is $1$, the
expression is nonnegative. If at least one is $1$, the second positive
floor is at least $1$. If both are $1$, then
$r+s+1\ge2q-2t+1\ge q$, so the first positive floor is also at least
$1$. This proves nonnegativity in all cases.
\end{proof}

\section*{The denominator bound and induction}
\begin{proof}[Proof of Theorem~\ref{main}]
Set $a=k-1$, $b=n-k-1$, and $t=\mu+1$. Expanding the binomial
coefficients gives the positive rational number
\begin{equation}\label{ratio}
 2C_\mu(n,k)=
 \frac{(t-1)!(a+b+1)!(a+b+t)!}
 {a!\,b!\,(a+t)!\,(b+t)!}.
\end{equation}
For a prime $\ell$, Legendre's formula gives
\[
 v_\ell\bigl(2C_\mu(n,k)\bigr)
 =\sum_{e\ge1} F_{\ell^e}(a,b,t).
\]
The sum is finite. By Lemma~\ref{floor}, all terms with
$\ell^e\ge2\mu+1$ are nonnegative, and each remaining term is at
least $-1$. Therefore
\[
 v_\ell\bigl(2C_\mu(n,k)\bigr)
 \ge-\#\{e\ge1:\ell^e\le2\mu\}
 =-v_\ell(L_\mu).
\]
For $\mu=0$ the set is empty and $L_0=1$. Thus every prime valuation
of $2L_\mu C_\mu(n,k)=D_\mu C_\mu(n,k)$ is nonnegative, proving
that this positive rational number is an integer.

Finally put $b_n=a_n/D_\mu$. The initial value is $b_1=1$, and
\eqref{rec} becomes
\[
 b_n=\sum_{k=1}^{n-1}
 \bigl(D_\mu C_\mu(n,k)\bigr)b_kb_{n-k}.
\]
Induction proves $b_n\in\mathbb Z_{>0}$ for every $n$, as required.
\end{proof}

\section*{Verification and source scope}
The proof is uniform in all indices and does not rely on finite testing.
The accompanying standard-library Python verifier checks 153,450 exact
recurrence coefficients, 9,300 factorial-ratio identities, 13,533,960
floor-function cases, 840 sequence/scaling identities, and 21 terms from
the source's examples. All pass. The whole current arXiv manuscript and
the 25-page journal article were inspected. The latter retains Conjecture
8.7 on page 22. Its postscript resolves the separate logconvexity
Conjecture 9.1 and does not assert a resolution of Conjecture 8.7.
This note was reviewed by its author agent; no independent peer review
is claimed. Dated status evidence and bounded literature-search limits
are recorded in the accompanying audit files.

\section*{Four follow-up problems}
\begin{enumerate}
\setcounter{enumi}{1}
\item \textbf{Task 2.} Determine the least positive integer $M_\mu$ such
that $M_\mu C_\mu(n,k)$ is integral for all $n\ge2$, $1\le k<n$,
and give its prime factorization as a function of $\mu$.
\item \textbf{Task 3.} Determine the least positive integer initial value
$A_\mu$ for which the entire sequence \eqref{rec} is integral. Compare
$A_\mu$ with $M_\mu$; cancellation within the recurrence is allowed.
\item \textbf{Task 4.} For the canonical initial value $a_1=2$, classify
the primes occurring in reduced denominators and determine
$\liminf_{n\to\infty}v_\ell(a_n)/n$ for each such prime $\ell$.
\item \textbf{Task 5.} Classify rational parameters $\mu>-1/2$ for which
a nonzero rational initial value makes every term of \eqref{rec} an
integer, using generalized binomial coefficients in that recurrence.
\end{enumerate}

\begin{thebibliography}{9}
\bibitem{AMV}
T. Amdeberhan, V. H. Moll and C. Vignat,
\emph{A probabilistic interpretation of a sequence related to Narayana
polynomials}, Online Journal of Analytic Combinatorics 8 (2013),
paper 1, 1--25.
\href{https://doi.org/10.61091/ojac-801}{doi:10.61091/ojac-801}.
Preprint \href{https://arxiv.org/abs/1202.1203}{arXiv:1202.1203}.
\bibitem{MathDB}
\href{https://mathdb.com/p/322992}{MathDB 322992},
\emph{The $p$-integrality conjecture for the Narayana-related sequence}.
\end{thebibliography}
\end{document}
